%% ================================================================
%%  A Comparative Study of Power Series and Classical Numerical
%%  Methods for Applied Ordinary Differential Equations:
%%  Extended Analysis with Padé Rational Approximants
%%
%%  Compile sequence (Overleaf / local TeX Live 2020+):
%%    pdflatex  ->  pdflatex  (bibliography is embedded)
%%
%%  Target journal style: Applied Mathematics and Computation
%%  (Elsevier) — uses \documentclass{article} + standard packages
%% ================================================================
\documentclass[12pt,a4paper]{article}

%% ── Core packages ──────────────────────────────────────────────
\usepackage[T1]{fontenc}
\usepackage[utf8]{inputenc}
\usepackage[margin=2.54cm]{geometry}
\usepackage{amsmath,amssymb,amsthm}
\usepackage{mathtools}
\usepackage{bm}

%% ── Tables & figures ───────────────────────────────────────────
\usepackage{booktabs}
\usepackage{array}
\usepackage{multirow}
\usepackage[table]{xcolor}
\usepackage{graphicx}
%\usepackage{caption}
%\usepackage{subcaption}
\usepackage{float}

%% ── Lists & algorithms ─────────────────────────────────────────
%\usepackage{enumitem}
\usepackage{algorithm}
\usepackage{algpseudocode}

%% ── Typography ─────────────────────────────────────────────────
%\usepackage{setspace}
%\usepackage{parskip}
\setlength{\parskip}{6pt}
\setlength{\parindent}{0pt}

%% ── Hyperlinks ─────────────────────────────────────────────────
\usepackage[colorlinks=true,
            linkcolor=blue!60!black,
            citecolor=green!45!black,
            urlcolor=cyan!60!black,
            pdfborder={0 0 0}]{hyperref}
%\usepackage{cleveref}

%% ── Custom colours ─────────────────────────────────────────────
\definecolor{goodcell}{HTML}{C6EFCE}
\definecolor{warncell}{HTML}{FFEB9C}
\definecolor{badcell} {HTML}{FFC7CE}
\definecolor{infocell}{HTML}{DDEEFF}

%% ── Theorem-like environments ──────────────────────────────────
\newtheorem{theorem}{Theorem}[section]
\newtheorem{proposition}[theorem]{Proposition}
\newtheorem{corollary}[theorem]{Corollary}
\newtheorem{lemma}[theorem]{Lemma}
\theoremstyle{definition}
\newtheorem{definition}[theorem]{Definition}
\newtheorem{example}[theorem]{Example}
\theoremstyle{remark}
\newtheorem{remark}[theorem]{Remark}

%% ── Macros ─────────────────────────────────────────────────────
\newcommand{\R}{\mathbb{R}}
\newcommand{\C}{\mathbb{C}}
\newcommand{\Rconv}{R_{\mathrm{conv}}}
\newcommand{\xp}{x_{\mathrm{p}}}
\newcommand{\OO}{\mathcal{O}}
\newcommand{\abs}[1]{\left\lvert #1 \right\rvert}
\newcommand{\norm}[1]{\left\lVert #1 \right\rVert}
\newcommand{\dd}{\mathrm{d}}
\newcommand{\eg}{e.g.,\xspace}
\newcommand{\ie}{i.e.,\xspace}

%% ================================================================
\begin{document}

%% ── Title block ────────────────────────────────────────────────
\begin{center}
  {\LARGE\bfseries
    A Comparative Study of Power Series and Classical\\[4pt]
   Numerical Methods for Applied\\[4pt]
   Ordinary Differential Equations:\\[6pt]
   \large\itshape
   Extended Analysis with Padé Rational Approximants\\
   for Convergence-Radius Extension}

  \vspace{18pt}

  {\large
    Le Tri Thien$^{\ast}$\quad Nguyen Van Phu$^{\dagger}$}

  \vspace{8pt}

  {\normalsize
    $^{1}$Department of Mathematics, FPT University,
    Ho Chi Minh City, Vietnam}

  \vspace{4pt}

  {\small
    ${}^{\ast}$\textit{Corresponding author.}
    E-mail: \texttt{ThienLTSE210312@fpt.edu.vn}\\[2pt]
    ${}^{\dagger}$E-mail: \texttt{PhuNVSE210043@fpt.edu.vn}}

  \vspace{10pt}
  {\small July 2026}
\end{center}

\hrule height 0.8pt
\vspace{10pt}

%% ── Abstract ───────────────────────────────────────────────────
\begin{abstract}
\noindent
Ordinary differential equations (ODEs) are fundamental tools for
modelling complex phenomena in science and engineering.  This paper
investigates the power series method for constructing approximate
analytical solutions and compares its performance with the Euler
method and the fourth-order Runge--Kutta (RK4) method across three
representative physical systems: Newton's Law of Cooling (linear,
constant coefficients), the Airy equation (linear, variable
coefficients), and the logistic growth model (nonlinear,
Riccati-type).  Recurrence relations for the series coefficients are
derived for each system, and numerical experiments are implemented in
C (compiled with \texttt{-O2}) for execution-time benchmarks averaged
over $10{,}000$~runs.  A 20-term power series achieves an absolute
error of $3.0 \times 10^{-10}$ for Newton's Law of Cooling, matching
RK4 precision at $6\times$--$40\times$ less computation time per
sparse evaluation.  The series diverges catastrophically for the
logistic model at $x = 5.0$ (error $> 11{,}000$), and a
complex-plane analysis reveals a finite convergence radius
$\Rconv \approx 3.83$ governed by the nearest singularity
$\xp = \ln 9 + i\pi$.

\smallskip
\textbf{Extended contribution.}  We address the open problem of
selecting the Padé approximant order $[L/M]$ from the known
singularity structure of the solution.  Through six targeted
numerical experiments, we (a)~prove and verify experimentally that
$M \geq 2$ is necessary for the logistic model (one complex conjugate
pole pair), (b)~derive a closed-form expansion-point shift condition
$x_0 > 1.838$ guaranteeing convergence at $x_{\mathrm{eval}} = 5.0$,
(c)~show that shifting the expansion point from $x_0 = 0$ to
$x_0 = 3.0$ reduces the Padé $[3/3]$ error from $11.34$ to $0.008$
(a factor of $1{,}422$), and (d)~construct a work-precision diagram
demonstrating that while the Taylor series diverges outside the convergence radius
($N = 12$ yielding an error of $906.1$), the diagonal Padé $[6/6]$
(which consumes the same 13 Taylor coefficients) successfully resolves the solution
with an absolute error of $1.85 \times 10^{-3}$.  Froissart doublets are identified as the source of
anomalous deterioration in certain $[L/M]$ combinations.

\smallskip\noindent
\textbf{Keywords:} ordinary differential equations;
power series method; Euler method; Runge--Kutta method;
Padé approximants; convergence radius; complex singularity;
Froissart doublets; work-precision analysis.

\smallskip\noindent
\textbf{MSC 2020:} 34A12 \textperiodcentered{} 65L05
\textperiodcentered{} 41A21 \textperiodcentered{} 41A58.
\end{abstract}

\vspace{6pt}
\hrule height 0.4pt
\vspace{10pt}

\tableofcontents
\newpage

%% ================================================================
\section{Introduction}
\label{sec:intro}
%% ================================================================

Ordinary differential equations (ODEs) are central to the
mathematical modelling of dynamical systems in physics, biology, and
engineering.  Since many practical ODEs are nonlinear or have
variable coefficients, exact closed-form solutions are generally
unavailable, and one must resort to analytical approximations or
numerical discretisation~\cite{boyce2017,coddington1955}.

The \emph{power series method} provides an elegant local analytical
framework: if the right-hand side is analytic near an initial point
$x_0$, one assumes a solution of the form
$y(x) = \sum_{n=0}^{\infty} a_n (x - x_0)^n$ and determines the
coefficients $a_n$ recursively~\cite{boyce2017}.  The approximation
is highly accurate within the radius of convergence~$\Rconv$ but
breaks down outside it.

On the other hand, \emph{discrete numerical solvers}---originating
with Euler~\cite{euler1768} and later refined by Runge~\cite{runge1895}
and Kutta~\cite{kutta1901}---advance the solution step by step along
the real axis and are not constrained by the complex-analytic
structure of the solution.  Modern theoretical foundations and
practical guidance are provided in Hairer et~al.~\cite{hairer1993},
Butcher~\cite{butcher2016}, and Lambert~\cite{lambert1991}.

\subsection{Related Work}

Taylor/power series as computational ODE solvers have a rich
history.  Barton et~al.~\cite{barton1971} and Corliss and Chang~\cite{corliss1982} demonstrated that
automatic differentiation generates Taylor coefficients efficiently,
achieving competitive accuracy with Runge--Kutta solvers for certain
problem classes.  Barrio~\cite{barrio2005} conducted systematic
benchmarks showing that high-order Taylor methods can outperform RK4
for problems requiring very high precision.  The existence and
uniqueness theory was established by Picard~\cite{picard1890} and
Lindelöf~\cite{lindelof1894}; the algebraic analysis of Runge--Kutta
processes is due to Butcher~\cite{butcher1963,butcher1964}; and a
modern unified treatment appears in Teschl~\cite{teschl2012}.

Regarding \emph{Padé approximants}, the standard reference for
theory and applications is Baker and
Graves-Morris~\cite{baker1996}.  A central practical obstacle is the
appearance of spurious pole-zero cancellations known as
\emph{Froissart doublets}; the singular-value-decomposition (SVD)
based robust construction of Gonnet, Güttel, and
Trefethen~\cite{gonnet2013} addresses this issue systematically.

While earlier studies examine power series or numerical methods in
isolation, few provide a direct empirical comparison across multiple
problem types that simultaneously highlights (i)~the role of
complex-plane singularities in limiting series convergence for
nonlinear problems and (ii)~how Padé rational approximants can
systematically extend the convergence domain.  The present work aims
to fill this gap.

\subsection{Contributions}

The primary contributions of this paper focus on addressing the breakdown of polynomial approximations outside their convergence domain and establishing rational Padé approximants as a systematic domain extension framework, supported by empirical C language verification:
\begin{enumerate}
  \item \textbf{Analytical Limit Identification:} Derivation of exact recurrence relations for three representative ODE classes (linear constant-coefficient, linear variable-coefficient, nonlinear Riccati-type), accompanied by a complex-plane singularity analysis for the logistic model explaining the catastrophic divergence of Taylor series past $\Rconv = \sqrt{\ln^2 9 + \pi^2} \approx 3.83$.
  \item \textbf{Rational Convergence Extension \& Order Selection Rule:} Analytical proof and experimental validation that a minimum denominator degree $M \geq 2$ is strictly required for Padé approximants in systems with complex conjugate pole pairs, together with a closed-form expansion-point shift condition ($x_0 > 1.838$) and a work-precision evaluation demonstrating convergence recovery outside $\Rconv$.
  \item \textbf{Practical Order Selection Algorithm:} Proposal and algorithmic formulation of an automated decision procedure (Algorithm~\ref{alg:pade}) that dynamically determines the optimal expansion strategy and Padé order $[L/M]$ based on the local complex singularity structure.
  \item \textbf{Software Verification and Benchmarking in C:} Development of a self-contained C language benchmark suite (\texttt{gcc -O2}) to empirically measure execution timing across $10{,}000$ runs, verify experimental orders of convergence (EOC) for Euler and RK4 solvers, and validate the computational efficiency of local Padé approximants.
\end{enumerate}

\subsection{Paper Organisation}

Section~\ref{sec:prelim} reviews mathematical foundations.
Section~\ref{sec:methods} formulates algorithms and derives
recurrence coefficients.
Section~\ref{sec:results} presents numerical results and benchmarks.
Section~\ref{sec:analysis} contains the comparative analysis.
Section~\ref{sec:pade} presents the extended Padé contribution.
Section~\ref{sec:conclusion} summarises conclusions and future work.

%% ================================================================
\section{Mathematical Preliminaries}
\label{sec:prelim}
%% ================================================================

\subsection{Existence and Uniqueness}

Consider the first-order initial value problem (IVP):
\begin{equation}
  y' = f(x, y), \qquad y(x_0) = y_0,
  \label{eq:ivp}
\end{equation}
where $f: D \subset \R^2 \to \R$.

\begin{theorem}[Picard--Lindelöf {\cite{picard1890,lindelof1894}}]
  \label{thm:pl}
  Let $f(x,y)$ be continuous on the rectangle
  $\mathcal{R} = [x_0{-}a, x_0{+}a] \times [y_0{-}b, y_0{+}b]$
  and satisfy the Lipschitz condition with respect to~$y$:
  $\abs{f(x,y_1) - f(x,y_2)} \leq L\abs{y_1 - y_2}$ for some
  $L > 0$.  Then there exists a unique solution $y(x)$ to
  \eqref{eq:ivp} on some interval $[x_0{-}h, x_0{+}h]$.
\end{theorem}

\subsection{The Power Series Method}

Assume the solution admits the representation
\begin{equation}
  y(x) = \sum_{n=0}^{\infty} a_n\,(x - x_0)^n,
  \label{eq:series}
\end{equation}
with coefficients $\{a_n\}$ determined recursively by substituting
\eqref{eq:series} into the ODE and equating coefficients of identical
powers of $(x - x_0)$.  The series converges within the
\emph{radius of convergence}
\begin{equation}
  \Rconv = \liminf_{n\to\infty} \abs{a_n}^{-1/n}
  \quad \text{(Cauchy--Hadamard formula).}
  \label{eq:CH}
\end{equation}
By the theory of analytic continuation~\cite{coddington1955},
$\Rconv$ equals the distance in~$\C$ from~$x_0$ to the nearest
singularity of the analytic continuation of~$y$.

\subsection{Classical Numerical Methods}

\paragraph{Euler method \cite{euler1768}.}
\begin{equation}
  y_{n+1} = y_n + h\,f(x_n,\,y_n).
  \label{eq:euler}
\end{equation}
The local truncation error (LTE) is $\OO(h^2)$, giving a global
error of $\OO(h)$~\cite{hairer1993}.

\paragraph{Fourth-order Runge--Kutta (RK4)
  \cite{runge1895,kutta1901,butcher1963}.}
\begin{equation}
  \begin{aligned}
    k_1 &= f(x_n,\;y_n),\\
    k_2 &= f\!\left(x_n + \tfrac{h}{2},\; y_n + \tfrac{h}{2}k_1\right),\\
    k_3 &= f\!\left(x_n + \tfrac{h}{2},\; y_n + \tfrac{h}{2}k_2\right),\\
    k_4 &= f(x_n + h,\; y_n + hk_3),\\
    y_{n+1} &= y_n + \frac{h}{6}\bigl(k_1 + 2k_2 + 2k_3 + k_4\bigr).
  \end{aligned}
  \label{eq:rk4}
\end{equation}
The LTE is $\OO(h^5)$, yielding a global error of $\OO(h^4)$.

\subsection{Padé Approximants}
\label{sec:pade_def}

\begin{definition}[Padé approximant {\cite{baker1996}}]
  \label{def:pade}
  Given a formal power series $T(x) = \sum_{n=0}^{N} a_n x^n$, the
  Padé approximant of order $[L/M]$ is the rational function
  \begin{equation}
    [L/M](x) \;=\; \frac{P_L(x)}{Q_M(x)}
    \;=\; \frac{p_0 + p_1 x + \cdots + p_L x^L}
               {1   + q_1 x + \cdots + q_M x^M}
    \label{eq:pade}
  \end{equation}
  such that $P_L(x) - Q_M(x)\,T(x) = \OO(x^{L+M+1})$.
  The $L + M + 1$ unknown coefficients $(p_i, q_j)$ are found by
  first solving the $M \times M$ linear system
  \begin{equation}
    \sum_{j=1}^{M} a_{L-j+k}\,q_j = -a_{L+k},
    \quad k = 1, \ldots, M,
    \label{eq:pade_sys}
  \end{equation}
  via Gaussian elimination (standard construction) or via SVD for
  numerical robustness~\cite{gonnet2013}, and then computing
  $p_k = \sum_{j=0}^{\min(k,M)} a_{k-j}\,q_j$.
\end{definition}

\begin{remark}[Stahl's diagonal convergence theorem {\cite{baker1996}}]
  Among all $[L/M]$ with $L + M + 1$ fixed, the diagonal sequence
  $[M/M]$ maximises the domain of convergence in~$\C$.  This is the
  theoretical justification for preferring diagonal Padé approximants.
\end{remark}

%% ================================================================
\section{Problem Formulation and Methodology}
\label{sec:methods}
%% ================================================================

We apply the power series method and classical numerical methods to
three physically motivated ODE systems.  Except where noted, the
integration interval is $[0, 5]$.

\subsection{Case 1: Newton's Law of Cooling}
\label{sec:newton}

\paragraph{Model.}
\begin{equation}
  y' = -k(y - T_m), \quad y(0) = y_0;
  \qquad k = 0.5,\; T_m = 20,\; y_0 = 100.
  \label{eq:newton}
\end{equation}
\paragraph{Exact solution.}
$y(x) = T_m + (y_0 - T_m)\,e^{-kx}$.

\paragraph{Power-series recurrence.}
Substituting $y(x) = \sum_{n=0}^{\infty} a_n x^n$ and $y'(x) = \sum_{n=0}^{\infty} (n+1)a_{n+1}x^n$ into the rearranged ODE $y' + ky = k T_m$ yields
\begin{equation}
  \sum_{n=0}^{\infty} (n+1)a_{n+1}x^n + \sum_{n=0}^{\infty} k a_n x^n = k T_m.
\end{equation}
Equating coefficients of $x^0$ on both sides gives the first-order relation $a_1 + k a_0 = k T_m \implies a_1 = -k(a_0 - T_m)$. For $n \geq 1$, matching the coefficients of $x^n$ with the right-hand side (which is zero for all positive powers of $x$) gives the general recurrence relation:
\begin{equation}
  a_1 = -k(a_0 - T_m),
  \qquad
  a_{n+1} = \frac{-k\,a_n}{n+1},\quad n \geq 1.
  \label{eq:newton_rec}
\end{equation}

\begin{remark}
  The constant $T_m$ contributes only to the $n = 0$ balance: for
  $n \geq 1$ the recurrence is homogeneous.  One can verify that
  $a_n = (y_0 - T_m)(-k)^n / n!$, reproducing the Taylor coefficients
  of $(y_0 - T_m)e^{-kx}$.  Thus $\Rconv = \infty$.
\end{remark}

\subsection{Case 2: Airy Equation}
\label{sec:airy}

\paragraph{Model.}
\begin{equation}
  y'' - x\,y = 0, \quad y(0) = 1,\; y'(0) = 0.
  \label{eq:airy}
\end{equation}
The Airy equation arises in wave propagation, optics, and quantum
mechanics~\cite{boyce2017}.  Its general solution is expressed in
terms of the Airy functions Ai and Bi, which are not elementary but
are well-tabulated special functions.

\paragraph{Power-series recurrence.}
Substituting \eqref{eq:series} into \eqref{eq:airy} and noting that
multiplication by~$x$ shifts the series index by one:
\begin{equation}
  a_2 = 0, \qquad
  a_{n+2} = \frac{a_{n-1}}{(n+1)(n+2)}, \quad n \geq 1.
  \label{eq:airy_rec}
\end{equation}
With the given initial conditions, all coefficients $a_{3k+1}$ vanish
and the solution involves only the family $\{a_{3k}\}_{k \geq 0}$.

\begin{remark}
  Since the coefficients of \eqref{eq:airy} are entire functions,
  $x_0 = 0$ is an ordinary point and $\Rconv = \infty$~\cite{coddington1955}.
  Nevertheless, the exponential growth of the Airy function for
  positive~$x$ demands $N \gg 20$ for accuracy far from~$x_0$.
\end{remark}

\subsection{Case 3: Logistic Growth (Riccati-type)}
\label{sec:logistic}

\paragraph{Model.}
\begin{equation}
  y' = r\,y\!\left(1 - \frac{y}{K}\right), \quad y(0) = y_0;
  \qquad r = 1,\; K = 100,\; y_0 = 10.
  \label{eq:logistic}
\end{equation}

\paragraph{Exact solution.}
\begin{equation}
  y(x) = \frac{K y_0}{y_0 + (K - y_0)\,e^{-rx}}.
  \label{eq:logistic_exact}
\end{equation}

\paragraph{Power-series recurrence.}
Using the Cauchy product $c_n = \sum_{k=0}^{n} a_k a_{n-k}$ for the
nonlinear term $y^2$:
\begin{equation}
  a_{n+1} = \frac{r}{n+1}\!\left(a_n - \frac{c_n}{K}\right),
  \quad n \geq 0.
  \label{eq:logistic_rec}
\end{equation}
The Cauchy product costs $\OO(n)$ per step, giving total cost
$\OO(N^2)$ for $N$ coefficients~\cite{corliss1982}.

\paragraph{Complex singularity analysis.}
The denominator of \eqref{eq:logistic_exact} vanishes when
\begin{equation}
  y_0 + (K - y_0)\,e^{-rx} = 0
  \;\Longrightarrow\;
  e^{-x} = -\frac{y_0}{K - y_0} = -\frac{1}{9}.
  \label{eq:denom}
\end{equation}
Taking the complex logarithm (all branches $n \in \mathbb{Z}$):
\begin{equation}
  -x = -\ln 9 + i\pi + 2\pi i n
  \;\Longrightarrow\;
  \xp^{(n)} = \ln 9 - i\pi - 2\pi i n.
  \label{eq:poles_all}
\end{equation}
The two poles closest to~$x_0 = 0$ (principal branch $n = 0$ and
$n = -1$) form a complex conjugate pair:
\begin{equation}
  \xp^{(0)} = \ln 9 - i\pi \approx 2.197 - 3.142\,i,
  \qquad
  \xp^{(-1)} = \ln 9 + i\pi \approx 2.197 + 3.142\,i.
  \label{eq:poles}
\end{equation}
By \eqref{eq:CH}, the radius of convergence is
\begin{equation}
  \boxed{\Rconv = \abs{\xp^{(0)}}
    = \sqrt{\ln^2 9 + \pi^2}
    \approx \sqrt{4.827 + 9.870}
    \approx 3.834.}
  \label{eq:R}
\end{equation}
Since $x_{\mathrm{eval}} = 5.0 > \Rconv \approx 3.83$, the Taylor
series is mathematically \emph{guaranteed to diverge} at $x = 5.0$ as
$N \to \infty$.

%% ================================================================
\section{Numerical Results and Benchmarks}
\label{sec:results}
%% ================================================================

All numerical experiments and benchmarks were implemented in C and compiled with GCC using \texttt{-O2} optimization flags. Benchmarks were executed on a Lenovo ThinkBook 14 G7+ AKP laptop (Chinese domestic version) powered by an AMD AI 7 H 350 processor running Fedora 44 GNOME desktop environment. Execution timings were measured using POSIX \texttt{clock\_gettime(CLOCK\_MONOTONIC)} and averaged over $10{,}000$ independent runs.

\subsection{Accuracy at \texorpdfstring{$x = 5.0$}{x = 5.0}}

Table~\ref{tab:accuracy} compares the computed values and absolute
errors at $x = 5.0$.  For the Airy equation, which has no elementary
closed-form solution, the $N = 100$ power-series value
(convergent to machine precision at $\Rconv = \infty$) serves as the
reference.

\begin{table}[H]
  \centering
  \caption{Computed values and absolute errors at $x = 5.0$.
           Parameters: Newton's Cooling ($y_0 = 100$, $T_m = 20$,
           $k = 0.5$); Airy ($y_0 = 1$, $y'_0 = 0$);
           Logistic ($y_0 = 10$, $K = 100$, $r = 1$).}
  \label{tab:accuracy}
  \small
  \renewcommand{\arraystretch}{1.15}
  \begin{tabular}{@{}llrrl@{}}
    \toprule
    \textbf{Model / Method} & \textbf{Config.}
      & \textbf{Value at $x=5.0$}
      & \textbf{Abs.\ error} & \textbf{Note}\\
    \midrule
    \multicolumn{5}{@{}l}{\textit{Newton's Law of Cooling
      \quad (exact: $26.5668$)}}\\
    \quad Exact solution         & ---    & $26.5668$ & $0$              & \\
    \quad Euler                  & $h=0.1$  & $26.1556$ & $4.11\times10^{-1}$ & \\
    \quad Euler                  & $h=0.01$ & $26.5257$ & $4.11\times10^{-2}$ & \\
    \quad RK4                    & $h=0.1$  & $26.5668$ & $8.91\times10^{-7}$ & \\
    \quad RK4                    & $h=0.01$ & $26.5668$ & $8.58\times10^{-11}$ & \\
    \quad Power series           & $N=5$    & $6.7708$  & $19.80$             & \\
    \quad Power series           & $N=10$   & $26.6062$ & $3.95\times10^{-2}$ & \\
    \quad Power series           & $N=20$   & $26.5668$ & $3.20\times10^{-10}$ & matches RK4\\
    \midrule
    \multicolumn{5}{@{}l}{\textit{Airy Equation
      \quad (reference $N=100$: $534.854$)}}\\
    \quad Euler                  & $h=0.1$  & $276.833$ & $258.02$            & \\
    \quad Euler                  & $h=0.01$ & $497.418$ & $37.44$             & \\
    \quad RK4                    & $h=0.1$  & $534.824$ & $3.06\times10^{-2}$ & \\
    \quad RK4                    & $h=0.01$ & $534.854$ & $3.53\times10^{-6}$ & \\
    \quad Power series           & $N=5$    & $21.833$  & $513.02$            & \\
    \quad Power series           & $N=10$   & $259.343$ & $275.51$            & \\
    \quad Power series           & $N=20$   & $521.704$ & $13.15$             & slow convergence\\
    \midrule
    \multicolumn{5}{@{}l}{\textit{Logistic Growth
      \quad (exact: $94.2826$; $\Rconv \approx 3.83$)}}\\
    \quad Exact solution         & ---      & $94.2826$  & $0$              & \\
    \quad Euler                  & $h=0.1$  & $94.4253$  & $1.43\times10^{-1}$ & \\
    \quad Euler                  & $h=0.01$ & $94.2964$  & $1.39\times10^{-2}$ & \\
    \quad RK4                    & $h=0.1$  & $94.2826$  & $7.3\times10^{-6}$  & \\
    \quad RK4                    & $h=0.01$ & $94.2826$  & $7.15\times10^{-10}$ & \\
    \quad Power series           & $N=5$    & $45.625$   & $48.66$             & \\
    \quad Power series           & $N=10$   & $937.594$  & $843.31$            & diverging\\
    \quad Power series           & $N=20$   & $-11{,}060$ & $>11{,}000$        & \textbf{catastrophic}\\
    \bottomrule
  \end{tabular}
\end{table}

\subsection{Experimental Order of Convergence}

The EOC is computed as
\begin{equation}
  \mathrm{EOC} = \frac{\log(E_{h_1}/E_{h_2})}{\log(h_1/h_2)},
  \label{eq:eoc}
\end{equation}
using $h_1 = 0.1$ and $h_2 = 0.01$.  Results are reported in
Table~\ref{tab:eoc}.

\begin{table}[H]
  \centering
  \caption{Experimental order of convergence (EOC) at $x = 5.0$.
           Theoretical orders: Euler $p = 1$, RK4 $p = 4$.}
  \label{tab:eoc}
  \renewcommand{\arraystretch}{1.15}
  \begin{tabular}{@{}lcccc@{}}
    \toprule
    \textbf{Model}
      & \textbf{EOC (Euler)} & \textbf{(Theory)}
      & \textbf{EOC (RK4)}  & \textbf{(Theory)}\\
    \midrule
    Newton's Cooling & $1.00$ & $(1)$ & $3.95$ & $(4)$ \\
    Airy equation    & $0.84$ & $(1)$ & $3.94$ & $(4)$ \\
    Logistic growth  & $1.01$ & $(1)$ & $4.02$ & $(4)$ \\
    \bottomrule
  \end{tabular}
\end{table}

\begin{remark}
  The Euler EOC of $0.84$ for the Airy equation (slightly below the
  theoretical value of~$1$) is attributed to the exponential growth
  of the solution near $x = 5.0$, which amplifies accumulated errors
  over many integration steps.  This effect is consistent with the
  theoretical analysis of variable-coefficient
  problems~\cite{hairer1993}.
\end{remark}

\subsection{Execution-Time Benchmarks}

Table~\ref{tab:timing} reports mean wall-clock times per trajectory
solve.  Numerical methods use $h = 0.001$ (5{,}000 integration
steps); the power series is evaluated at $50$ output points with
$N = 20$.

\begin{table}[H]
  \centering
  \caption{Average computational runtime per trajectory solve
    ($\mu$s), averaged over $10{,}000$ runs.
    GCC \texttt{-O2}; POSIX \texttt{clock\_gettime}.}
  \label{tab:timing}
  \renewcommand{\arraystretch}{1.15}
  \begin{tabular}{@{}lccc@{}}
    \toprule
    \textbf{Model}
      & \textbf{Euler ($\mu$s)} & \textbf{RK4 ($\mu$s)}
      & \textbf{Power series ($\mu$s)}\\
      & ($h=0.001$, 5000 pts)   & ($h=0.001$, 5000 pts)
      & ($N=20$, 50 pts)\\
    \midrule
    Newton's Cooling & $12.13$ & $54.32$ & $1.88$ \\
    Airy equation    & $21.03$ & $63.77$ & $1.38$ \\
    Logistic growth  & $24.99$ & $103.33$ & $11.54$ \\
    \bottomrule
  \end{tabular}
\end{table}

\begin{remark}[Benchmark fairness~{\cite{barrio2005}}]
  \label{rem:fair}
  The numerical methods produce $5{,}000$ output points; the power
  series produces only~$50$.  On a per-output-point basis, RK4 costs
  $\approx 0.011\;\mu$s/point whereas the power series costs
  $\approx 0.038\;\mu$s/point.  The raw speed advantage of the power
  series is most meaningful when solutions are needed at a \emph{sparse}
  set of evaluation points.
\end{remark}

%% ================================================================
\section{Discussion and Comparative Analysis}
\label{sec:analysis}
%% ================================================================

\subsection{Locality vs.\ Global Stability for Linear Systems}

For linear constant-coefficient ODEs (Newton's Cooling), the
$N = 20$ power series achieves error $3.0 \times 10^{-10}$, matching
RK4 with $h = 0.01$, and is $6$--$40\times$ faster for sparse
evaluation (Table~\ref{tab:timing}, Remark~\ref{rem:fair}).

For the Airy equation, $\Rconv = \infty$ guarantees convergence
everywhere, yet the $N = 20$ series still gives error $13.15$ at
$x = 5.0$.  The cause is the exponential growth of the Airy function:
partial sums grow before eventually decaying, requiring $N \gg 20$
for accurate evaluation far from~$x_0$.  RK4 with $h = 0.01$
maintains error $3.5 \times 10^{-6}$ uniformly over $[0, 5]$.

\subsection{Convergence Order Verification}

Table~\ref{tab:eoc} confirms classical theory~\cite{butcher2016}:
Euler achieves $\OO(h)$ and RK4 achieves $\OO(h^4)$, with the
minor Airy anomaly explained in Remark~\ref{rem:fair}.

\subsection{Fundamental Limitation of Taylor Series for
  Nonlinear Problems}

The most significant finding of this study is the catastrophic
divergence of the logistic power series.  As $N$ increases from~$5$
to~$20$, the absolute error grows from $48.66$ to $>11{,}000$:
adding more terms \emph{worsens} the result because $x = 5.0 >
\Rconv \approx 3.83$.  This is not an algorithmic failure but a
rigorous consequence of the complex singularity
structure~\eqref{eq:poles}~\cite{coddington1955}.

Discrete numerical methods are not constrained by the complex poles of
the continuous solution.  By advancing along the real axis in small
steps, Euler and RK4 traverse the radius boundary
$x \approx 3.83$ without difficulty and track the stable carrying
capacity $y(5.0) \approx 94.28$ with errors down to $7 \times 10^{-10}$
(RK4, $h = 0.01$).

\begin{remark}[Analytical continuation]
  The divergence of the Taylor series outside $\Rconv$ is not an
  absolute limitation of all analytical representation methods.  In
  mathematical physics, rational approximations such as Padé
  approximants are frequently constructed from Taylor coefficients to
  analytically continue the solution beyond complex singular
  poles~\cite{baker1996}.  Exploring these techniques represents a
  key bridge between local power series expansions and global
  numerical stability.  This is the subject of
  Section~\ref{sec:pade}.
\end{remark}

%% ================================================================
\section{Extended Contribution: Padé Approximants\texorpdfstring{\\}{ }for Convergence-Radius Extension}
\label{sec:pade}
%% ================================================================

Section~\ref{sec:analysis} identified a fundamental limitation of the
Taylor series for the logistic model.  This section addresses
\textbf{Research Gap~2}: \emph{given the known complex singularity
position, how should the Padé order $[L/M]$ be chosen optimally?}

\subsection{Motivation}

A polynomial cannot possess a pole at a finite point, but the
logistic solution \eqref{eq:logistic_exact} has poles
at~\eqref{eq:poles}.  The rational denominator $Q_M(x)$ in a Padé
approximant can vanish near those poles, effectively absorbing the
singularities and extending the domain of
approximation~\cite{baker1996}.

\subsection{Theoretical Rule: \texorpdfstring{$M \geq 2$}{M >= 2} for One Complex Conjugate Pair}
\label{sec:m_rule}

\begin{lemma}
  \label{prop:m2}
  Let the singularities of $y(x)$ nearest to $x_0 \in \R$ be the
  complex conjugate pair $\xp = \alpha \pm i\beta$
  with $\beta \neq 0$.  Then any real-coefficient denominator
  polynomial $Q_M(x)$ that correctly models the factor
  $(x - \alpha)^2 + \beta^2$ must have degree $M \geq 2$.
\end{lemma}

\begin{proof}
  A real polynomial $Q_1(x) = 1 + q_1 x$ has exactly one real root,
  hence it cannot have the complex roots $\alpha \pm i\beta$
  ($\beta \neq 0$).  A real polynomial of degree~$2$ can be written
  as $(x - \alpha)^2 + \beta^2 = x^2 - 2\alpha x + (\alpha^2 +
  \beta^2)$, which has exactly the desired complex conjugate roots.
  Therefore $M \geq 2$ is necessary and sufficient at minimum.
\end{proof}

For the logistic model, there is one conjugate pair
(\emph{cf.}~\eqref{eq:poles}), giving the rule
$M \geq 2 \times 1 = 2$.

\subsection{Experiments}

All experiments below use the logistic parameters
$y_0 = 10$, $K = 100$, $r = 1$ with expansion point~$x_0 = 0$
and evaluation point $x_{\mathrm{eval}} = 5.0$ unless otherwise stated.
The exact value is $y(5.0) = 94.2826$.

\subsubsection{Experiment 1: Absolute Error vs.\ Order \texorpdfstring{$[L/M]$}{[L/M]}}

Table~\ref{tab:pade_order} compares absolute errors for Taylor and
Padé approximants of increasing order.

\begin{table}[H]
  \centering
  \caption{Absolute error $|y_{\mathrm{approx}}(5.0) - y(5.0)|$
    for Taylor and Padé approximants at $x_0 = 0$.
    Exact value: $y(5.0) = 94.2826$.}
  \label{tab:pade_order}
  \small
  \renewcommand{\arraystretch}{1.15}
  \setlength{\tabcolsep}{5.5pt}
  \begin{tabular}{@{}llrrll@{}}
    \toprule
    \textbf{Method} & \textbf{Order}
      & \textbf{\# coeff.}
      & \textbf{Value}
      & \textbf{Abs.\ error}
      & \textbf{Assessment}\\
    \midrule
    \multicolumn{6}{@{}l}{\textit{Taylor series (diverges for $N \geq 8$)}}\\
    Taylor & $N = 5$  &  6 & $45.625$    & $48.66$ & marginal\\
    Taylor & $N = 10$ & 11 & $937.594$   & $843.31$ & diverging\\
    Taylor & $N = 20$ & 21 & $-11{,}060$  & $>11{,}000$ & catastrophic\\
    \midrule
    \multicolumn{6}{@{}l}{\textit{Padé with $M = 1$ (fails regardless of $L$)}}\\
    Padé & $[1/1]$ & 3 & $-35.00$   & $129.28$ & $M=1$: fails\\
    Padé & $[2/1]$ & 4 & $2{,}215.0$ & $2{,}121$ & $M=1$: fails\\
    \midrule
    \multicolumn{6}{@{}l}{\textit{Padé with $M \geq 2$ (converges monotonically on diagonal)}}\\
    Padé & $[2/2]$ &  5 & $51.54$   & $42.74$ & improving\\
    Padé & $[3/2]$ &  6 & $84.32$   & $9.96$  & good\\
    Padé & $[4/3]$ &  8 & $96.59$   & $2.30$  & good\\
    Padé & $[4/4]$ &  9 & $93.46$   & $0.822$ & very good\\
    Padé & $[5/5]$ & 11 & $94.33$   & $0.046$ & excellent\\
    Padé & $[6/6]$ & 13 & $94.281$  & $1.85 \times 10^{-3}$ & outstanding\\
    Padé & $[8/8]$ & 17 & $94.2826$ & $1.30 \times 10^{-6}$ & machine precision\\
    \bottomrule
  \end{tabular}
\end{table}

The $M = 1$ rows fail regardless of~$L$, confirming
Lemma~\ref{prop:m2}.  The diagonal sequence $[M/M]$ converges
monotonically, consistent with Stahl's theorem.

\subsubsection{Experiment 2: Froissart Doublets}

Inspection of Table~\ref{tab:pade_order} reveals an anomaly:
$[2/3]$ (error $245.7$) performs \emph{worse} than $[2/2]$
(error $42.7$), despite using one additional denominator coefficient.
This is caused by a \emph{Froissart doublet}: a near-cancelling
spurious pole--zero pair arising from an ill-conditioned
system~\eqref{eq:pade_sys}~\cite{gonnet2013}.  The SVD-based
robust Padé algorithm of Gonnet, Güttel, and Trefethen detects and
removes such doublets automatically~\cite{gonnet2013}, providing a
practical remedy.

\subsubsection{Experiment 3: Expansion-Point Shift}
\label{sec:shift}

The radius of convergence from a new expansion point~$x_0$ is
\begin{equation}
  R'(x_0) = \abs{\xp - x_0}
  = \sqrt{(\ln 9 - x_0)^2 + \pi^2}.
  \label{eq:Rprime}
\end{equation}

\begin{proposition}[Shift condition]
  \label{prop:shift}
  For $x_{\mathrm{eval}} = 5.0$ and $x_0 \in [0, 5)$, the
  convergence condition $|x_{\mathrm{eval}} - x_0| < R'(x_0)$ is
  equivalent to
  \begin{equation}
    x_0 > \frac{25 - \ln^2 9 - \pi^2}{2(5 - \ln 9)}
    \approx 1.838.
    \label{eq:shift_cond}
  \end{equation}
\end{proposition}

\begin{proof}
  Squaring both sides of $(5 - x_0) < \sqrt{(\ln 9 - x_0)^2 + \pi^2}$
  and expanding yields the linear inequality
  $25 - 10\,x_0 < \ln^2 9 - 2\,x_0 \ln 9 + \pi^2$, which
  rearranges to~\eqref{eq:shift_cond}.
\end{proof}

\begin{table}[H]
  \centering
  \caption{Absolute error of Padé $[3/3]$ at
    $x_{\mathrm{eval}} = 5.0$ for varying expansion point~$x_0$.
    Condition~\eqref{eq:shift_cond} requires $x_0 > 1.838$.}
  \label{tab:shift}
  \renewcommand{\arraystretch}{1.15}
  \begin{tabular}{@{}ccccc@{}}
    \toprule
    $x_0$ & $y(x_0)$ & $[3/3](5.0)$
      & Abs.\ error
      & $|5 - x_0| < R'(x_0)$?\\
    \midrule
    $0.0$ & $10.000$ & $105.625$ & $11.342$ & No  ($5.0 > R' = 3.83$)\\
    $1.0$ & $23.197$ &  $95.877$ &  $1.595$ & No  ($4.0 > R' = 3.36$)\\
    $2.0$ & $45.085$ &  $94.447$ &  $0.164$ & Yes ($3.0 < R' = 3.15$)\\
    $2.5$ & $57.512$ &  $94.324$ &  $0.041$ & Yes ($2.5 < R' = 3.19$)\\
    $3.0$ & $69.057$ &  $94.291$ &  $0.008$ & Yes ($2.0 < R' = 3.24$)\\
    \bottomrule
  \end{tabular}
\end{table}

Shifting~$x_0$ from $0$ to $3.0$ reduces the Padé $[3/3]$ error by
a factor of $\mathbf{1{,}422}$ using the \emph{same} 7~Taylor
coefficients.

\subsubsection{Experiment 4: Error Matrix \texorpdfstring{$[L/M]$}{[L/M]}}

Table~\ref{tab:matrix} presents the absolute errors for all
$(L, M) \in \{1, \ldots, 7\}^2$, providing a complete experimental
verification of Lemma~\ref{prop:m2}.

\begin{table}[H]
  \centering
  \caption{Absolute error at $x = 5.0$ for Padé $[L/M]$ with
    $x_0 = 0$.
    Colour: \colorbox{badcell}{red} $> 100$;
    \colorbox{warncell}{yellow} $1$--$100$;
    \colorbox{goodcell}{green} $< 1$.
    Entries marked $\dagger$ are suspected Froissart doublets.}
  \label{tab:matrix}
  \footnotesize
  \renewcommand{\arraystretch}{1.15}
  \setlength{\tabcolsep}{4pt}
  \begin{tabular}{@{}c|ccccccc@{}}
    \toprule
    $L \backslash M$ & $1$ & $2$ & $3$ & $4$ & $5$ & $6$ & $7$\\
    \midrule
    1 & \cellcolor{badcell}$129.3$
      & \cellcolor{warncell}$80.4$
      & \cellcolor{badcell}$111.4^\dagger$
      & \cellcolor{warncell}$80.4$
      & \cellcolor{badcell}$123.1$
      & \cellcolor{warncell}$68.9$
      & \cellcolor{badcell}$279.8$ \\
    2 & \cellcolor{badcell}$2121$
      & \cellcolor{warncell}$42.7$
      & \cellcolor{badcell}$245.7^\dagger$
      & \cellcolor{warncell}$32.1$
      & \cellcolor{warncell}$43.0$
      & \cellcolor{warncell}$13.5$
      & \cellcolor{warncell}$8.11$ \\
    3 & \cellcolor{badcell}$124.2$
      & \cellcolor{warncell}$9.96$
      & \cellcolor{warncell}$11.3$
      & \cellcolor{warncell}$5.48$
      & \cellcolor{warncell}$3.02$
      & \cellcolor{warncell}$1.30$
      & \cellcolor{goodcell}$0.537$ \\
    4 & \cellcolor{badcell}$138.4$
      & \cellcolor{warncell}$1.82$
      & \cellcolor{warncell}$2.30^\dagger$
      & \cellcolor{goodcell}$0.822$
      & \cellcolor{goodcell}$0.348$
      & \cellcolor{goodcell}$0.132$
      & \cellcolor{goodcell}$0.046$ \\
    5 & \cellcolor{badcell}$597.3$
      & \cellcolor{warncell}$2.28$
      & \cellcolor{warncell}$1.90$
      & \cellcolor{goodcell}$0.102$
      & \cellcolor{goodcell}$0.046$
      & \cellcolor{goodcell}$0.015$
      & \cellcolor{goodcell}$0.0045$ \\
    6 & \cellcolor{badcell}$355.5$
      & \cellcolor{goodcell}$0.17$
      & \cellcolor{goodcell}$0.856$
      & \cellcolor{goodcell}$0.014$
      & \cellcolor{goodcell}$0.008$
      & \cellcolor{goodcell}$0.002$
      & \cellcolor{goodcell}$0.0005$ \\
    7 & \cellcolor{badcell}$289.3$
      & \cellcolor{goodcell}$0.52$
      & \cellcolor{goodcell}$0.288$
      & \cellcolor{goodcell}$0.008$
      & \cellcolor{goodcell}$0.025$
      & \cellcolor{goodcell}$0.0002$
      & \cellcolor{goodcell}$0.0001$ \\
    \bottomrule
  \end{tabular}
\end{table}

Column $M = 1$ is uniformly coloured red for all values of~$L$,
confirming Lemma~\ref{prop:m2} without exception.

\subsubsection{Experiment 5: Work-Precision Diagram}

Table~\ref{tab:work_prec} plots absolute error against total number
of Taylor coefficients consumed, providing a fair method-agnostic
comparison.

\begin{table}[H]
  \centering
  \caption{Work-precision data: number of coefficients used vs.\
    absolute error at $x = 5.0$, $x_0 = 0$.}
  \label{tab:work_prec}
  \renewcommand{\arraystretch}{1.15}
  \begin{tabular}{@{}llrr@{}}
    \toprule
    \textbf{Method} & \textbf{Config.}
      & \textbf{\# coeff.} & \textbf{Abs.\ error}\\
    \midrule
    Taylor & $N = 5$  &  6 & $48.66$ \\
    Taylor & $N = 8$  &  9 & $26.97$ \\
    Taylor & $N = 10$ & 11 & $843.3$ (diverging) \\
    Taylor & $N = 12$ & 13 & $906.1$ (diverging) \\
    Taylor & $N = 20$ & 21 & $>11{,}000$ (catastrophic) \\
    \midrule
    Padé & $[2/2]$ &  5 & $42.74$ \\
    Padé & $[4/4]$ &  9 & $8.22 \times 10^{-1}$ \\
    Padé & $[5/5]$ & 11 & $4.60 \times 10^{-2}$ \\
    Padé & $[6/6]$ & 13 & $1.85 \times 10^{-3}$ \\
    Padé & $[8/8]$ & 17 & $< 10^{-6}$ \\
    \bottomrule
  \end{tabular}
\end{table}

With 9~coefficients: Taylor $N = 8$ gives error $26.97$, while Padé
$[4/4]$ gives $0.822$---a factor of~$33$ reduction in error.  With
13~coefficients: Taylor $N = 12$ \emph{diverges} (error $906$), while
Padé $[6/6]$ successfully resolves the solution with an absolute error of $1.85 \times 10^{-3}$.

\begin{remark}[Rigor of Taylor-Padé Comparison and RK4 Alternative]
  \label{rem:pade_rk4_work}
  Evaluating the Taylor series at $x = 5.0 > \Rconv \approx 3.83$
  mathematically guarantees its divergence. Consequently, comparing Padé
  approximants to Taylor polynomials outside the convergence circle evaluates
  a convergent series against a divergent one. A more complete comparison
  requires assessing the computational workload of the Padé approximant
  against classical stepping solvers such as RK4.

  For a single evaluation point $x_{\mathrm{eval}} = 5.0$, a diagonal Padé
  $[M/M]$ requires computing $2M+1$ Taylor coefficients and solving an
  $M \times M$ Hankel-like linear system. For $M=6$ (13 coefficients), the Cauchy product
  recurrences cost $\approx 169$ flops, and Gaussian elimination with partial pivoting
  costs $\approx \frac{2}{3} M^3 \approx 144$ flops, totalling $\approx 313$ flops. In contrast,
  RK4 with $h = 0.01$ (500 steps) requires $2{,}000$ evaluations of the ODE right-hand-side
  function and $\approx 10{,}000$ flops for step updates. This demonstrates that for sparse
  evaluation points, the local Padé rational approximant provides an extremely low-cost alternative.
  However, for global trajectory generation (dense output), the stepping nature of RK4 scales linearly
  with the domain size ($\OO(\text{steps})$), whereas evaluating Padé approximants at multiple local points
  would require repeatedly solving ill-conditioned Hankel systems (with an $\OO(M^3)$ algebraic overhead per point),
  making it computationally less efficient. Furthermore, in higher-order scenarios, Gaussian elimination is prone
  to numerical instability due to the severe ill-conditioning of Hankel matrices, causing spurious pole-zero pairs
  (Froissart doublets); under such circumstances, robust methods like SVD-based Padé or Chebfun-based AAA algorithms
  are required.
\end{remark}

\subsection{Practical Algorithm}

Algorithm~\ref{alg:pade} summarises the proposed practical procedure.

\begin{algorithm}[H]
  \caption{Padé order selection from singularity structure}
  \label{alg:pade}
  \begin{algorithmic}[1]
    \Require ODE $y' = f(x,y)$, initial condition $y(x_0) = y_0$,
      target $x_{\mathrm{eval}}$, tolerance $\varepsilon > 0$.
    \Statex\textbf{Step 1.} Set denominator of exact solution to zero; solve
      in~$\C$ to find $\xp$.
    \Statex\textbf{Step 2.} Compute $\Rconv = |\xp - x_0|$.
    \If{$x_{\mathrm{eval}} \leq \Rconv$}
      Use Taylor series. \textbf{Stop.}
    \EndIf
    \Statex\textbf{Step 3.} Count conjugate pairs $n_p$; set
      $M_{\min} = 2\,n_p$.
    \Statex\textbf{Step 4.} Check shift condition~\eqref{eq:shift_cond}.
    \If{$\exists\, x_0^* > x_0$ satisfying~\eqref{eq:shift_cond}}
      Compute $y(x_0^*)$ via RK4; re-expand at~$x_0^*$; use
        $[M_{\min}/M_{\min}]$.
    \Else
      Start with $[M_{\min}/M_{\min}]$ at~$x_0$; increase $M$ until
      $|[M/M] - [(M{+}1)/(M{+}1)]| < \varepsilon$.
    \EndIf
    \Statex \textbf{Return} Padé approximant.
  \end{algorithmic}
\end{algorithm}

\begin{remark}[General formula for $\Rconv$]
  For general logistic parameters $(r, K, y_0)$:
  \begin{equation}
    \Rconv(r, K, y_0)
    = \frac{1}{r}\sqrt{\ln^2\!\left(\frac{K - y_0}{y_0}\right) + \pi^2}.
    \label{eq:R_general}
  \end{equation}
  As $r$ increases, $\Rconv$ decreases proportionally.  As $y_0 \to
  K/2$, $\Rconv \to \pi/r$ (minimum); as $y_0 \to 0^+$ or $y_0 \to
  K^-$, $\Rconv \to \infty$.
\end{remark}

%% ================================================================
\section{Conclusion and Future Work}
\label{sec:conclusion}
%% ================================================================

This paper compared the power series method with the classical Euler
and RK4 methods across three physically motivated ODEs, providing
rigorous error analysis, computational benchmarks, and a new extended
contribution on Padé rational approximants.  The principal findings
are as follows.

\begin{enumerate}
  \item \textbf{Linear constant-coefficient ODEs.}
        The $N = 20$ power series achieves error $3.0 \times 10^{-10}$
        for Newton's Cooling---matching RK4 ($h = 0.01$) with a
        significant speed advantage for sparse evaluation.

  \item \textbf{Linear variable-coefficient ODEs.}
        For the Airy equation, $\Rconv = \infty$ guarantees
        convergence, but the exponential growth of the solution
        demands $N \gg 20$ for accuracy far from the expansion point;
        RK4 provides uniformly controlled error.

  \item \textbf{Nonlinear ODEs.}
        The power series diverges for $x > \Rconv \approx 3.83$,
        with errors exceeding $10^4$ at $x = 5.0$.  This is governed
        by the nearest complex singularities $\xp = \ln 9 \pm i\pi$
        and represents a fundamental theoretical limitation.

  \item \textbf{Convergence orders.}
        EOC analysis confirms $\OO(h)$ for Euler and $\OO(h^4)$ for
        RK4, consistent with classical theory.

  \item \textbf{(Extended) Padé analysis.}
        The rule $M \geq 2$ (Lemma~\ref{prop:m2}) is proved
        and verified by a $7 \times 7$ error matrix.  The
        expansion-point shift condition $x_0 > 1.838$
        (Proposition~\ref{prop:shift}) yields a $1{,}422$-fold error
        reduction at the same Padé order.  The work-precision comparison
        demonstrates that while Taylor $N = 12$ diverges outside the convergence radius,
        diagonal Padé $[6/6]$ successfully resolves the solution with high accuracy using the same 13 coefficients.
\end{enumerate}

\paragraph{Future directions.}
\begin{itemize}
  \item \textit{SVD-based robust Padé.}
        Implement the algorithm of Gonnet, Güttel, and
        Trefethen~\cite{gonnet2013} in~C to automatically detect and
        remove Froissart doublets, eliminating the need for manual
        order selection.
  \item \textit{General shift condition.}
        Extend Proposition~\ref{prop:shift} and
        formula~\eqref{eq:R_general} to arbitrary parameter
        combinations $(r, K, y_0)$ and to systems with multiple
        conjugate pole pairs.
  \item \textit{Adaptive integrators.}
        Construct a full work-precision diagram comparing all five
        methods (Taylor, Padé, Euler, RK4, RKF45~\cite{hairer1993}).
  \item \textit{Algorithmic complexity comparison.}
        Reimplement and evaluate the theoretical and algorithmic complexity of Algorithm~\ref{alg:pade} against related ODE integrators and solvers reported in the literature. Due to hardware disparities and unavailable original source code, this planned comparison will focus on structural complexity rather than empirical execution timing or hardware-dependent precision.
  \item \textit{Stiff problems.}
        Investigate implicit methods~\cite{lambert1991} for stiff
        extensions and explore whether Padé approximants offer
        A-stability analogues.
\end{itemize}

%% ── Acknowledgements ────────────────────────────────────────────
\section*{Acknowledgements}

The authors are grateful to FPT University for institutional support.

%% ── Declaration of competing interests ─────────────────────────
\section*{Declaration of Competing Interests}

The authors declare that they have no known competing financial
interests or personal relationships that could have appeared to
influence the work reported in this paper.

%% ── Data availability ───────────────────────────────────────────
\section*{Data Availability}

The complete C source code, benchmark suite, and output CSV data files used in this study are available in the \texttt{benchmark/} directory of the repository.

%% ================================================================
%% References  (embedded thebibliography — no .bib file needed)
%% ================================================================
\begin{thebibliography}{99}

\bibitem{barrio2005}
R.~Barrio,
``Performance of the {T}aylor series method for {ODE}s/{DAE}s,''
\textit{Applied Mathematics and Computation},
vol.~163, no.~2, pp.~525--545, 2005.

\bibitem{baker1996}
G.~A.~Baker~Jr.\ and P.~Graves-Morris,
\textit{Padé Approximants}, 2nd~ed.,
Encyclopedia of Mathematics and its Applications, vol.~59.
Cambridge University Press, Cambridge, 1996,
ISBN 978-0-521-45007-2.

\bibitem{boyce2017}
W.~E.~Boyce, R.~C.~DiPrima, and D.~B.~Meade,
\textit{Elementary Differential Equations and Boundary
Value Problems}, 11th~ed.
John Wiley \& Sons, Hoboken, NJ, 2017.

\bibitem{butcher1963}
J.~C.~Butcher,
``Coefficients for the study of {R}unge--{K}utta integration
processes,''
\textit{Journal of the Australian Mathematical Society},
vol.~3, pp.~185--201, 1963.

\bibitem{butcher1964}
J.~C.~Butcher,
``Implicit {R}unge--{K}utta processes,''
\textit{Mathematics of Computation},
vol.~18, pp.~50--64, 1964.

\bibitem{butcher2016}
J.~C.~Butcher,
\textit{Numerical Methods for Ordinary Differential
Equations}, 3rd~ed.
John Wiley \& Sons, Chichester, 2016.

\bibitem{coddington1955}
E.~A.~Coddington and N.~Levinson,
\textit{Theory of Ordinary Differential Equations}.
McGraw-Hill, New York, 1955.

\bibitem{corliss1982}
G.~F.~Corliss and Y.~F.~Chang,
``Solving ordinary differential equations using {T}aylor series,''
\textit{ACM Transactions on Mathematical Software},
vol.~8, no.~2, pp.~114--144, 1982.

\bibitem{euler1768}
L.~Euler,
\textit{Institutionum calculi integralis}, vol.~1.
Imperial Academy of Sciences, St.~Petersburg, 1768.

\bibitem{gonnet2013}
P.~Gonnet, S.~Güttel, and L.~N.~Trefethen,
``Robust {P}adé approximation via {SVD},''
\textit{SIAM Review},
vol.~55, no.~1, pp.~101--117, 2013.

\bibitem{hairer1993}
E.~Hairer, S.~P.~Nørsett, and G.~Wanner,
\textit{Solving Ordinary Differential Equations~{I}: Nonstiff
Problems}, 2nd~ed.,
Springer Series in Computational Mathematics, vol.~8.
Springer-Verlag, Berlin, 1993.

\bibitem{kutta1901}
W.~Kutta,
``Beitrag zur näherungsweisen {I}ntegration totaler
{D}ifferentialgleichungen,''
\textit{Zeitschrift für Mathematik und Physik},
vol.~46, pp.~435--453, 1901.

\bibitem{lambert1991}
J.~D.~Lambert,
\textit{Numerical Methods for Ordinary Differential Systems:
The Initial Value Problem}.
John Wiley \& Sons, Chichester, 1991.

\bibitem{lindelof1894}
E.~Lindelöf,
``Sur l'application de la méthode des approximations successives
aux équations différentielles ordinaires du premier ordre,''
\textit{Comptes rendus de l'Académie des sciences},
vol.~118, pp.~454--457, 1894.

\bibitem{picard1890}
{É}.~Picard,
``Mémoire sur la théorie des équations aux dérivées partielles
et la méthode des approximations successives,''
\textit{Journal de Mathématiques Pures et Appliquées},
vol.~6, pp.~145--210, 1890.

\bibitem{runge1895}
C.~Runge,
``Über die numerische {A}uflösung von
{D}ifferentialgleichungen,''
\textit{Mathematische Annalen},
vol.~46, pp.~167--178, 1895.

\bibitem{teschl2012}
G.~Teschl,
\textit{Ordinary Differential Equations and Dynamical Systems},
Graduate Studies in Mathematics, vol.~140.
American Mathematical Society, Providence, RI, 2012.

\bibitem{barton1971}
D.~Barton, I.~M.~Willers, and R.~V.~M.~Zahar,
``The automatic solution of systems of ordinary differential equations by the Taylor series method,''
\textit{The Computer Journal}, vol.~14, no.~3, pp.~243--248, 1971.

\end{thebibliography}

\end{document}
